\chapter{Conclusion and Limitations}

The whole pipeline can be reproduced from start to end. DVC pins the data version, every run
records its commit, data version, seed, parameters, metrics and artifacts in MLflow on DagsHub, the
served model is picked by a written rule on validation data and promoted with a registry alias, and
the system starts with one \cmd{docker compose up}. The model API runs on Azure behind HTTPS and an
API key, and the web app stores every decision with its time, model version and confidence.

The model itself has clear limits. The dataset is small and scraped. The AI images come from a few
generators, and the ``real'' images include edited and stock photos, classical paintings and
clip-art. The error analysis and Grad-CAM (\cref{sec:2-6}) show that the model partly learned
style, treating painterly images as AI and photographic ones as real, instead of real generation
artefacts. On new generators, screenshots or images re-compressed by social media it will do worse
than the test numbers suggest. Size, format and compression also differ between the classes
(\cref{ch:data}), so even after resizing some of what the model learned may be those shortcuts.

The metrics are noisy as well. A single train/validation/test split of a few hundred images moves
noticeably with small changes, and cross-validation or a bigger, more varied dataset would make the
comparison in \cref{sec:3-5} more reliable. On the deployment side, one VM is a single point of
failure, and there is no autoscaling, rate limiting or monitoring of prediction drift yet.

If I continued the project, I would train on a larger dataset with more generators, add JPEG
re-compression as a robustness test at evaluation time, log production predictions back to MLflow
to watch for drift, and have CI deploy the API image on every tagged release.
